% AzDIM-Bench — paper skeleton (arXiv preprint style)
% Build: pdflatex main && bibtex main && pdflatex main && pdflatex main
\documentclass[11pt]{article}
\usepackage[margin=1in]{geometry}
\usepackage{booktabs,graphicx,amsmath,url,xcolor}
\usepackage[hidelinks]{hyperref}
\usepackage[numbers]{natbib}

\title{AzDIM-Bench: A Contamination-Aware Benchmark from\\
Azerbaijan's National Examinations for Evaluating LLMs}
\author{Javid Mardanov\\
\small with AI-assisted dataset construction and evaluation tooling}
\date{2026}

\begin{document}
\maketitle

\begin{abstract}
% ~180 words. Numbers from results/metrics_v0.json and results/dim_scores_v0.json.
Azerbaijani, spoken by roughly thirty million people, is nearly absent from
LLM evaluation. We introduce AzDIM-Bench, a benchmark of 2{,}249 items
built from the official published archives of Azerbaijan's State Examination
Center (D\.IM): 11th-grade school-leaving examinations and two-stage
university-entrance examinations (specialty groups I--IV) from 2025--2026,
spanning 14 subjects and six languages (Azerbaijani and Russian sectors plus
foreign-language sections). Items are transcribed from official explanation
PDFs by a schema-constrained vision model, cross-checked deterministically
against D\.IM's separately published answer-key grids, and human-verified.
We evaluate 14 frontier and open-weight models under a reason-then-answer
protocol, report accuracy with the official 700-point admission-score
simulation (including stage-2 negative marking), and exploit the exams'
public calendar for a natural post-training-cutoff split.
% Key findings sentence: language gap AZ<RU, literature hardest, small-model
% negative-marking collapse, 12.5->100 format effect.
\end{abstract}

\section{Introduction}
% Motivation: low-resource eval gap; DIM's unusual publication practice as a
% measurement instrument; contributions list (dataset, protocol, QC method,
% score simulation, findings).

\section{Background: the Azerbaijani examination system}
% From docs/research/exam_system_research.md; include the two-stage figure
% (export site scorecanvas SVG to PDF for \includegraphics).

\section{Related work}
% From docs/research/related_work.md (agent survey).

\section{Dataset construction}
\subsection{Sources and provenance}      % manifest, SHA-256, izah vs etalon
\subsection{Vision-based extraction}     % text-layer corruption evidence; 3-page window; structured outputs
\subsection{Canonicalization}            % fragments, variant dedup, subjects
\subsection{Quality control}             % deterministic etalon cross-check; concordant-position agreement; human verification protocol
\subsection{The extractor-circularity caveat}  % mitigations

\section{Evaluation protocol}
% Tracks A/B/F; reason-then-answer rationale with the 12.5%->100% experiment;
% hidden-reasoning minimization; parsing; invalid handling; core-set design.

\section{Results}
% Leaderboard table (results/metrics_v0.json); language-gap dumbbell;
% 2025-vs-2026 split; subject breakdown; invalid rates.

\section{D\.IM score simulation}
% Official formula (cite DIM PDF); central vs floor assumptions;
% 700-point table per group; negative-marking collapse of small models.

\section{Contamination analysis}
% 2025 vs 2026 deltas per model; discussion of what it can and cannot show.

\section{Limitations}
% written/multimodal items excluded from main track; extraction residual error
% bounds from QC; single-run at fixed condition; Grade-9 question text
% unavailable publicly.

\section{Release}
% HF dataset, code, site, license position, DIM attribution.

\bibliographystyle{unsrtnat}
\bibliography{refs}
\end{document}
